\documentclass[10pt,a4paper]{amsart}
\usepackage[margin=19mm]{geometry}
\usepackage{amsmath,amssymb,mathtools}
\usepackage[scr=rsfs,cal=euler]{mathalfa}
\usepackage{xfrac}
\usepackage[unicode,hidelinks,bookmarks=false]{hyperref}
\newcommand{\ie}{i.e.\ }
\newcommand{\cf}{cf.\ }
\newcommand{\resp}{respectively\ }
\newcommand{\Subsection}{\subsection}
\providecommand{\nobreakdash}{\nobreak\hbox{-}}
\setlength{\emergencystretch}{2em}
\multlinegap0pt
\newtheorem{lemma}{Lemma}
\newcommand{\bbZ}{\mathbb Z}
\title{Some remarks on L-equivalence for cubic fourfolds and hyper-Kähler manifolds}
\author{Simone Billi, Lucas Li Bassi}
\date{Source: arXiv:2512.01968v3 (CC BY 4.0)}
\begin{document}
\maketitle

\noindent\emph{Source and licence.} Some remarks on L-equivalence for cubic fourfolds and hyper-Kähler manifolds. Version arXiv:2512.01968v3 (CC BY 4.0), distributed by arXiv under the Creative Commons Attribution 4.0 International licence, http://creativecommons.org/licenses/by/4.0/, as recorded in the arXiv licence metadata for this exact version and shown on the abstract page https://arxiv.org/abs/2512.01968v3. Full licence notice: \texttt{licenses/arxiv-2512.01968-CC-BY-4.0.txt}.\par\smallskip

\noindent\textit{Editorial scope.} Counted unit: the article's lemma lemma:diofanteo (source0.tex lines 420--422). The article's complete original proof of it is delivered with it (source0.tex lines 423--434, one \texttt{proof} environment of 12 lines). No mathematics is written, repaired or re-derived: the statement and the proof are the article's own lines, in the article's own notation, and no retained line is substituted or re-flowed.  No further source-local context is needed: every cross-reference of every retained line resolves inside the two delivered spans.  Nothing was omitted from any retained span, so the package carries no omission record.  No retained line contains a citation, so the wrapper carries no bibliography and no bibliography entry.  No retained line contains a figure environment, an image-inclusion command or drawing code, so no image is delivered and the package declares no asset dependency.  Editorial formatting only, in addition: an AMS \texttt{amsart} preamble that restates the article's own declarations and macros used by the retained lines, namely the environments \texttt{lemma} on separate counters, together with the standard packages those lines need.  The retained environments print in this document as Lemma 1 in order of appearance, whereas the article prints the same items with the numbering of its own full document.  The complete native source of the article as distributed in the arXiv e-print for this exact version is bundled unchanged as \texttt{sources/190171/source0.tex}.
\begin{lemma}\label{lemma:diofanteo}
    Let $v, w\in A_2$ be two primitive elements with the same square. Then $v^{\perp}\simeq w^{\perp}$ as lattices.
\end{lemma}

\begin{proof}
    Let $v=ae_1+be_2$ for a basis $\{e_1, e_2\}$ of $A_2$. Then a generator of $v^{\perp}$ will be a primitive element $u=xe_1+ye_2$ such that $(2a-b)x+(2b-a)y=0$. We call $p:=(2a-b)$ and $q:=(2b-a)$. The diophantine equation $px+qy=0$ has solutions \begin{equation*}
        x=k\frac{q}{g},\ y=-k\frac{p}{g}, \qquad g:=\gcd(p,q),\quad k\in\bbZ.
    \end{equation*}
    The only primitive solutions are thus given by $k=\pm 1$. In both cases we get \begin{equation*}
        u^2=\frac{2}{g^2}\left(q^2+p^2+pq\right)=\frac{6}{g^2}\left(a^2+b^2-ab\right)=\frac{3}{g^2}v^2.
    \end{equation*}
    Note that $p\equiv q\equiv 2(a+b) \mod 3$. Moreover, $g\mid (2p+q)=3a$ and $g\mid (2q+p)=3b$. As $v$ is primitive we know that $\gcd(a,b)=1$, therefore $g\mid 3$, i.e. $g\in\left\{\pm 1,\pm 3\right\}$. Now, we have the identity $a^2+b^2-ab\equiv(a+b)^2 \mod 3$. Therefore,\begin{equation*}
        0\equiv v^2=2(a^2+b^2-ab) \mod 3\iff 0\equiv a+b\equiv p \equiv q\mod 3\iff g=\pm 3.
    \end{equation*}
    This shows that the square of the primitive element $u$ that generates the orthogonal complement of a primitive vector $v\in A_2$ is $u^2=\frac{v^2}{3}$ if $v^2\equiv 0\mod 3$ or, else, $u^2=3v^2$. In both cases it is uniquely determined by the square of $v$. 
\end{proof}

\end{document}
